\chapter*{Abstract}
\addcontentsline{toc}{chapter}{Abstract}

A shared--memory multiprocessor must answer a question that a uniprocessor never
has to face: when one core writes to memory, when --- and in what order relative
to its other writes --- must the other cores see that write?
The answer a machine
gives is called its \emph{memory consistency model}.
It is not an implementation
detail but an architectural contract, as binding as the instruction set itself,
and it is the reason a multithreaded C program that is provably correct on paper
can fail on real silicon.

The obvious contract is \emph{sequential consistency} (SC), formalised by Lamport
in 1979: the machine must behave as though the memory operations of all cores were
interleaved into a single total order that respects each core's program order.
SC matches programmer intuition exactly.
By the early 1990s, however, the
architecture community had rejected it for hardware implementation on performance
grounds, and every commercially significant instruction set since --- SPARC, Alpha,
POWER, ARM, and RISC-V --- has published a weaker contract, exposing reordering to
software and supplying fence instructions with which programmers are expected to
suppress it selectively.

That decision was reached on machines with in-order, non-speculative pipelines,
small bus-snooped caches, and core counts in the single digits.
Every one of those
premises has since been invalidated. A modern out-of-order core already contains
both mechanisms that a strong model needs: a coherence protocol that announces
remote writes, and a squash-and-restart path built for branch misprediction
recovery.
A line of work beginning with Gniady et al.\ has argued that, given
these, SC can be enforced speculatively at a small fraction of its assumed cost.
And yet no shipping processor implements it.

This paper examines that gap. It develops a single formal framework in which SC,
TSO, PSO, weak ordering and release consistency appear as successively weaker
constraints on one set of ordering relations; it traces the historical argument
from 1979 to the ratification of RVWMO; it establishes, by direct measurement on
x86 and ARM hardware, which relaxations are observable in practice as opposed to
merely permitted; and it quantifies, in a configurable multicore simulator, the
cost of enforcing each model with and without speculative load execution.
The
governing question is whether the performance objection to sequential consistency
still holds on a machine that already speculates --- and, if it does not, what
else explains the industry's continued preference for weak ordering.

\vspace{1em}
\noindent\textbf{Note on this version.} This is draft v0.1. It comprises the
formal, historical and methodological portions of the paper (Chapters~1--6).
The experimental and simulation results announced in Chapter~6 are the subject of
versions v0.2 and v0.3; Chapter~6 states the design of those experiments in
sufficient detail to be evaluated in advance of the data.
